\section{Introduction}\label{sec:intro}

Data marketplaces built on confidential computing promise that a data owner can let a buyer run
an analysis over a private dataset without handing the dataset to anyone: the data is decrypted
only inside a hardware-protected confidential virtual machine (CVM), and the platform operator
sees ciphertext~\cite{nascimento2020dnat,tokuda2026vddpi,xiao2020privacyguard}. This protects
the data \emph{from the operator}. It does not protect the data \emph{from the application}.
An application that computes the promised statistics can also hide individual records in the
low-order digits of a mean, in a base64 field, or in a group too small to be anonymous, and the
platform will dutifully deliver that output to the buyer.

The problem is sharpest in a common case: \emph{the code is known and the data is private}. When
the code is secret and the data is public, nobody is at risk. When the code is public, it can in
principle be audited, but auditing every application that every buyer submits is prohibitively
expensive, and the data owner is left with no practical guarantee. Using an LLM as the auditor is
tempting, yet LLMs alone detect vulnerabilities close to chance on rigorous
benchmarks~\cite{ding2025primevul}, and an adversarial application can address the reviewer
directly through comments and strings.

This paper asks a single question:
\emph{can a low-cost judge, armed with mechanical evidence, detect leaks hidden by expensive attacker models?}
We answer it with an \emph{AI Judge} that combines cheap, language-agnostic dynamic tests with an
LLM that reasons over their results, and we evaluate it in an attacker-versus-defender battle in
the style of AI control evaluations~\cite{greenblatt2024aicontrol}.

\textbf{Contributions.}
\begin{itemize}
  \item \textbf{C1: a layered AI Judge} (\S\ref{sec:judge}) that decides whether an application
  may run on a dataset given the owner's usage tags. Its core is a relational dynamic layer
  (canaries, one-record swaps and secret variation with black-box leakage
  estimation~\cite{cherubin2019fbleau}) whose findings become evidence for an LLM, followed by an
  output gate at real execution time.
  \item \textbf{C2: the Judge as a platform application} (\S\ref{sec:design}). The Judge runs
  through the same confidential pipeline as any other application and receives another
  application as input. Model, prompt and version are pinned by hash, so verdicts are
  reproducible and the Judge can be upgraded without changing the platform.
  \item \textbf{C3: a hidden-leakage benchmark} (\S\ref{sec:benchmark}) of aggregation,
  transformation and machine-learning tasks with benign and leaking variants and a
  capture-the-flag ground truth: an attack counts only if the task is solved \emph{and} a
  decoder recovers the canary from the output alone. To our knowledge no such benchmark exists.
  \item \textbf{C4: confidential execution of arbitrary OCI images} (\S\ref{sec:design}): attested
  two-phase builds and single-use flat CVMs that only talk to the platform, so that the
  legitimate output is the only remaining channel and the Judge's job is well defined.
\end{itemize}

\todo{parágrafo de resultados: detecção do defensor barato com e sem evidências, falsos positivos,
custo por veredito, custo do executor OCI em CVM.}

The rest of the paper is organized as follows. \S\ref{sec:background} reviews the building
blocks, \S\ref{sec:threat} defines the threat model, \S\ref{sec:design} presents the platform,
\S\ref{sec:judge} the Judge, \S\ref{sec:benchmark} the benchmark and \S\ref{sec:eval} the
evaluation. \S\ref{sec:related} and \S\ref{sec:discussion} cover related work and limitations.
